% Part: lambda-calculus
% Chapter: lambda-definability
% Section: lists

\documentclass[../../../include/open-logic-section]{subfiles}

\begin{document}

\olfileid{lam}{ldf}{lst}
\olsection{তালিকা}

$[M_1, M_2, \ldots, M_n]$ তালিকাটি এভাবে সংজ্ঞায়িত করা যায়:
\[
  [M_1, M_2, \ldots, M_n] \ident \lambd[sz][s M_1 (s M_2 (\ldots (s M_n z)))]
\]
অনানুষ্ঠানিকভাবে, তালিকাকে উপস্থাপনকারী পদটি তালিকার একটি
``সঞ্চায়ক'' অপেক্ষক: এটি দুটি পদ গ্রহণ করে। প্রথমটি এমন একটি অপেক্ষক,
যা সঞ্চিত মান ও একটি নতুন উপাদান গ্রহণ করে নতুন সঞ্চিত মান ফেরত দেয়;
দ্বিতীয়টি হলো প্রাথমিক সঞ্চিত মান। এভাবে উপাদানগুলি একে একে সঞ্চয়
করে এটি শেষে ফল ফেরত দেয়।

\begin{digress}
  অপেক্ষকভিত্তিক প্রোগ্রামিংয়ের সঙ্গে পরিচিত হলে এই সংজ্ঞাটির সঙ্গে
  \emph{ডান দিক থেকে ভাঁজ করা} বা right fold-এর মিল দেখতে পাবে:
  তালিকাকে তার উপাদানগুলির উপর ডান দিক থেকে ভাঁজ করার অপেক্ষক হিসেবে
  সংজ্ঞায়িত করা হয়েছে।
\end{digress}

এই সংকেতায়নের ভিত্তিতে কয়েকটি উপযোগী অপেক্ষক সহজেই সংজ্ঞায়িত করা যায়:
\begin{align*}
  \fn{Sum} & \ident
  \lambd[l][l \, (\lambd[xa][\fn{Add} \, x \, a])\,  \num{0}]\\
  \fn{Len} & \ident
  \lambd[l][l \, (\lambd[xa][\fn{Add} \, a \, \num{1}]) \num{0}]
\end{align*}

$\fn{Sum}$ চার্চ সংখ্যাপদের একটি তালিকার যোগফল নির্ণয় করে। এতে
তালিকাটির উপর সঞ্চয় চালানো হয়: প্রাথমিক মান $\num{0}$, আর প্রতিটি
উপাদানের জন্য নতুন মান হয় $\fn{Add} \, x\, a$ (এখানে $x$ বর্তমান
উপাদান এবং $a$ সঞ্চিত মান)। ফল হলো সব উপাদানের যোগফল।

\begin{prob}
  $\fn{Len}$ কী করে? ব্যাখ্যা করো।
\end{prob}

\begin{prob}
  একটি তালিকার ক্রম বিপরীত করে এমন একটি অপেক্ষক সংজ্ঞায়িত করো।
\end{prob}
\end{document}
